\section{แนวคิดของงานวิจัยและวิธีการดำเนินงาน}
\label{sec:method}

\subsection{แนวคิดของงานวิจัย}

% แนวคิดหลักของงานวิจัยนี้คือการแยกการควบคุมบทสนทนาออกจากการตีความคำตอบ ส่วนควบคุมบทสนทนาทำงานแบบกำหนดแน่นอนตามแบบประเมิน
% เพื่อให้ผู้ป่วยทุกคนได้รับคำถามที่มีถ้อยคำเหมือนกันและตรงตามต้นฉบับ ขณะที่โมเดลภาษาทำหน้าที่เฉพาะการจับคู่คำตอบภาษาพูดของผู้ป่วยเข้ากับระดับคะแนนของข้อคำถาม
% ระบบที่นำเสนอประกอบด้วย 6 องค์ประกอบ ดังรูปที่~\ref{fig:system-arch} และทุกองค์ประกอบทำงานบนเซิร์ฟเวอร์ภายในโรงพยาบาล
% โดยไม่มีการส่งข้อมูลเสียงหรือข้อความของผู้ป่วยไปยังบริการภายนอก

แนวคิดหลักของงานวิจัยนี้คือการใช้ส่วนควบคุมบทสนทนาแบบ finite-state machine ในการควบคุมการสนทนากับผู้ป่วยตามแบบประเมินที่กำหนดไว้ล่วงหน้า 
โดยให้ผู้ป่วยตอบคำถามทีละข้อผ่านเสียงพูด และใช้ LLMs ในการตีความคำตอบของผู้ป่วยเพื่อจับคู่กับตัวเลือกที่กำหนดไว้ในแบบประเมิน 
หลังจากผู้ป่วยตอบคำถามครบทุกข้อ ระบบจะคำนวณคะแนนรวมตามสูตรคำนวณของแบบประเมินที่ใช้และแสดงผลลัพธ์ให้ผู้ป่วยทราบ
ซึ่งระบบที่นำเสนอประกอบด้วย 6 ส่วน ดังรูปที่~\ref{fig:system-arch} และทุกองค์ประกอบทำงานบนเซิร์ฟเวอร์ภายในโรงพยาบาล
โดยไม่มีการส่งข้อมูลเสียงหรือข้อความของผู้ป่วยไปยังบริการภายนอก

\begin{figure}[htbp]
\centering
\includegraphics[width=0.65\textwidth]{figures/system_architecture.png}
\caption{สถาปัตยกรรมของระบบที่นำเสนอ}
\label{fig:system-arch}
\end{figure}

\begin{enumerate}
    \item \textbf{ePRO Form (JSON Format)}: คือ ePRO ที่ถูกแปลงให้อยู่ในรูปแบบ JSON เรียบร้อยแล้ว ซึ่งโครงสร้างจะประกอบด้วยข้อคำถาม 
    ระดับคำตอบเชิงอันดับพร้อมคำบรรยาย ไฟล์เสียงและภาพประกอบของแต่ละข้อ และประเภทของแบบประเมิน (DASH หรือ TESS)
    เพื่อให้สามารถนำไปงานต่อไปอย่างสะดวกและทำให้เพิ่มแบบประเมินใหม่ได้โดยไม่ต้องแก้ไขโปรแกรม
    \item \textbf{Dialogue Manager}: ควบคุมบทสนทนาแบบ finite-state machine 
    ที่จะเป็นตัวกำหนดลำดับการถามคำถาม และรับผลค่าความมั่นใจจาก option extractor 
    หากค่าความมั่นใจของการจับคู่สูงกว่าเกณฑ์ ระบบจะบันทึกตัวเลือกและไปยังข้อถัดไป
    หากต่ำกว่าเกณฑ์ ระบบจะเล่นเสียงขอให้ตอบใหม่หรือเสียงอ่านตัวเลือกที่สร้างไว้ล่วงหน้า
    และเมื่อถามซ้ำครบจำนวนครั้งที่กำหนดแล้วยังไม่ได้คำตอบที่ชัดเจน จะบันทึกข้อนั้นไว้ให้พยาบาลตรวจสอบ
    \item \textbf{Question Presenter}: เล่นเสียงคำถามที่ถูกสังเคราะห์ด้วย OmniVoice \citep{Zhu2026OmniVoice} และตรวจสอบไว้ล่วงหน้าแล้วพร้อมแสดงภาพประกอบ เพื่อให้ถ้อยคำคงที่และช่วยให้ผู้ป่วยเข้าใจคำถาม
    \item \textbf{Speech Recognizer}: บันทึกเสียงคำตอบของผู้ป่วยทีละข้อ แล้วแปลงเป็นข้อความด้วยโมเดล Typhoon Whisper Turbo 
    จากนั้นส่งข้อความไปยัง option extractor เพื่อจับคู่กับตัวเลือกของข้อคำถาม
    \item \textbf{Option Extractor}: ทำการจับคู่ข้อความคำตอบเข้ากับระดับคะแนนของข้อคำถามด้วย Qwen3.5 โดยใช้
    prompt ที่ระบุข้อคำถาม ตัวเลือกแต่ละระดับพร้อมคำบรรยาย และคำตอบของผู้ป่วย
    แล้วให้โมเดลสร้างผลลัพธ์เป็น JSON ตาม schema ที่กำหนด ซึ่งประกอบด้วย 4 ส่วน ได้แก่
    เหตุผลสั้น ๆ ที่อ้างอิงคำในคำตอบของผู้ป่วย ระดับคะแนนที่เลือก
    เจตนาของผู้ป่วยในกรณีที่ไม่ได้ตอบคำถาม เช่น ขอให้ทวนคำถามหรือไม่เข้าใจคำถาม และค่าความมั่นใจ
    ในระหว่างการสร้างผลลัพธ์จะใช้ constrained decoding บังคับให้ผลลัพธ์เป็นไปตาม schema
    และจำกัดระดับคะแนนให้อยู่ในชุดค่าของข้อคำถามนั้นเท่านั้น ส่งผลให้ผลลัพธ์อยู่ในรูปแบบที่โปรแกรมอ่านได้เสมอและไม่มีค่าที่อยู่นอกมาตราส่วน
    จากนั้นจึงส่งระดับคะแนน ค่าความมั่นใจ และเจตนาของผู้ป่วยให้ dialogue manager ดำเนินการต่อไป
    \item \textbf{Score Calculator}: ทำการแปลงตัวเลือกรายข้อเป็นคะแนนและคำนวณคะแนนรวมตามสูตรของแบบประเมินแต่ละฉบับด้วยโปรแกรม
    เพื่อให้การคิดคะแนนเป็นไปตามแบบประเมินต้นฉบับและทำสามารถทำซ้ำได้
\end{enumerate}

\subsection{วิธีการดำเนินงาน}

\begin{enumerate}
    \item ศึกษาทฤษฎีและงานวิจัยที่เกี่ยวข้องกับ PROMs ตัวแทนสนทนาในทางการแพทย์ การรู้จำเสียงพูดภาษาไทย และการตีความคำตอบด้วยโมเดลภาษา
    \item จัดทำแบบประเมิน DASH และ TESS ในรูปแบบโครงสร้างข้อมูล JSON และจัดทำเสียงคำถาม เสียงขอให้ตอบใหม่ เสียงอ่านตัวเลือก และภาพประกอบ
    \item ออกแบบและพัฒนา dialogue manager แบบ finite-state machine และเชื่อมต่อกับ question presenter และ speech recognizer
    \item จัดทำชุดข้อมูลคำตอบสังเคราะห์ภาษาไทยสำหรับการทดสอบ option extractor
    \item พัฒนา option extractor ด้วย Qwen3.5 เลือกขนาดของโมเดล และทดสอบกับชุดข้อมูลสังเคราะห์เพิ่อวัดความถูกต้องและประสิทธิภาพของโมเดล
    \item บันทึกเสียงการสัมภาษณ์โดยพยาบาล และประเมิน option extractor กับคำตอบจริงของผู้ป่วย
    \item integrate ระบบทั้งหมดเข้าด้วยกันและทดสอบการทำงานของระบบ
    \item นำระบบไปใช้กับผู้ป่วยจริงและวัดเวลาที่พยาบาลใช้
    \item วิเคราะห์ผล สรุปผลการวิจัย และจัดทำวิทยานิพนธ์
\end{enumerate}

\subsection{แผนการดำเนินงาน}

\definecolor{plan}{HTML}{62A69B}
\newcommand{\g}{\cellcolor{plan}}   % คำสั่งย่อสำหรับระบายสีช่อง
\newcolumntype{C}[1]{>{\centering\arraybackslash}m{#1}}

\begin{table}[htbp]
  \centering
  \renewcommand{\arraystretch}{1.6}
  \setlength{\tabcolsep}{4pt}
  \begin{tabular}{|C{6.2cm}|*{7}{C{1.05cm}|}}
    \hline
    \textbf{แผนการดำเนินงาน} & \textbf{ส.ค.} & \textbf{ก.ย.} & \textbf{ต.ค.} & \textbf{พ.ย.} & \textbf{ธ.ค.} & \textbf{ม.ค.} & \textbf{ก.พ.} \\
    \hline
    ศึกษาทฤษฎีและงานวิจัยที่เกี่ยวข้อง
      & \g & \g &    &    &    &    &    \\ \hline
    จัดทำแบบประเมิน DASH และ TESS ในรูปแบบ JSON รวมถึงไฟล์อื่นที่เกี่ยวข้อง
      &    & \g &    &    &    &    &    \\ \hline
    ออกแบบและพัฒนา dialogue manager
      &    & \g & \g &    &    &    &    \\ \hline
    จัดทำชุดข้อมูลคำตอบสังเคราะห์ภาษาไทย
      &    &    & \g &    &    &    &    \\ \hline
    พัฒนา option extractor และทดสอบกับข้อมูลสังเคราะห์
      &    &    & \g & \g &    &    &    \\ \hline
    บันทึกเสียงการสัมภาษณ์ และประเมิน option extractor
    &    &    &    & \g & \g &    &    \\ \hline
    integrate ระบบทั้งหมดเข้าด้วยกันและทดสอบ
      &    &    &    &  & \g &    &    \\ \hline
    ทดลองใช้ระบบกับผู้ป่วยจริงควบคู่กับการสัมภาษณ์โดยพยาบาล
      &    &    &    &    & \g & \g &    \\ \hline
    วิเคราะห์ผล สรุปผลการวิจัย และจัดทำวิทยานิพนธ์
      &    &    &    &    &    & \g & \g \\ \hline
  \end{tabular}
  \caption{แผนการดำเนินงาน}
\label{tab:plan}
\end{table}

\newpage
\subsection{การทดลองและการประเมินผล}

\subsubsection{การทดลองที่ 1: การเลือกโมเดลสำหรับ Option Extractor}

ทดสอบ option extractor กับชุดทดสอบของชุดข้อมูลคำตอบสังเคราะห์ โดยเปรียบเทียบ Qwen3.5 ขนาดต่าง ๆ ที่มีขนาดไม่เกิน 9 พันล้านพารามิเตอร์
ซึ่งเป็นขนาดที่ทำงานได้บนเซิร์ฟเวอร์ภายใน ทุกขนาดใช้ prompt และ schema เดียวกัน ปิดโหมดคิด และกำหนดค่าการสร้างข้อความให้คงที่
การทดลองนี้วัดผล 2 ด้าน ดังนี้
\begin{enumerate}
    \item \textbf{ความถูกต้องและความสอดคล้องของการจับคู่} เทียบระดับคะแนนที่โมเดลเลือกกับระดับคะแนนเป้าหมายของข้อมูลสังเคราะห์
    โดยวัดสัดส่วนของข้อที่เลือกระดับคะแนนตรงกันด้วย accuracy 
    และวัด weighted kappa เพื่อประเมินความสอดคล้องของการจับคู่ที่คำนึงถึงระยะห่างของระดับคะแนน
    \item \textbf{ทรัพยากรการคำนวณ} วัดเวลาที่ใช้ตีความคำตอบหนึ่งข้อ ภายใต้หน่วยความจำของหน่วยประมวลผลกราฟิก (GPU) ที่มีอยู่อย่างจำกัด
\end{enumerate}
ผลการทดลองนี้ใช้เลือกขนาดของโมเดลที่สมดุลระหว่างความถูกต้องกับทรัพยากรการคำนวณ เพื่อนำไปใช้ในการทดลองที่ 2 และ 3

\subsubsection{การทดลองที่ 2: ความสอดคล้องกับการสัมภาษณ์โดยพยาบาล}

การทดลองนี้ประเมิน option extractor กับคำตอบจริงของผู้ป่วย โดยบันทึกเสียงระหว่างการสัมภาษณ์ตามแบบประเมินโดยพยาบาลตามขั้นตอนปกติ
ทั้งเสียงคำถามของพยาบาลและเสียงคำตอบของผู้ป่วย และไม่บันทึกข้อมูลที่ระบุตัวตนของผู้ป่วย
โดยที่จะใช้คะแนนรายข้อที่พยาบาลบันทึกระหว่างการสัมภาษณ์เป็นคะแนนอ้างอิง รวมถึงจับระยะเวลาที่ใช้ในการสัมภาษณ์ด้วย
ซึ่งจะเก็บข้อมูลจากผู้ป่วยเป็นจำนวน 40 คน แบ่งเป็น 20 คนสำหรับแบบประเมิน DASH และ 20 คนสำหรับแบบประเมิน TESS

เสียงที่บันทึกจะถูกแบ่งเป็นช่วงตามข้อคำถาม โดยคำตอบของแต่ละข้อคือเสียงของผู้ป่วยตั้งแต่พยาบาลถามข้อนั้นจนถึงก่อนพยาบาลถามข้อถัดไป
จากนั้นแปลงเป็นข้อความด้วย speech recognizer และจับคู่กับระดับคะแนนด้วย option extractor
แล้ววัดความสอดคล้องกับคะแนนของพยาบาลในระดับรายข้อด้วย weighted kappa
และในระดับคะแนนรวมด้วยค่า ICC โดยคำนวณแยกตามแบบประเมินแต่ละฉบับ

\subsubsection{การทดลองที่ 3: การใช้งานระบบกับผู้ป่วยจริง}

การทดลองนี้นำระบบทั้งหมดไปใช้เก็บข้อมูลกับผู้ป่วยจริง เพื่อประเมินว่าระบบช่วยลดเวลาที่พยาบาลต้องใช้ในการเก็บข้อมูลได้หรือไม่
โดยที่จะเก็บข้อมูลจากผู้ป่วยเท่ากับการทดลองที่ 2
เวลาที่บันทึกมี 2 ค่า ได้แก่
\begin{enumerate}
    \item \textbf{เวลาที่ผู้ป่วยใช้กับระบบ} บันทึกโดยอัตโนมัติตั้งแต่เริ่มเล่นเสียงคำถามข้อแรกจนได้คำตอบของข้อสุดท้าย
    \item \textbf{เวลาที่พยาบาลใช้ตรวจสอบ} บันทึกตั้งแต่พยาบาลเริ่มตรวจสอบข้อที่ระบบส่งให้ตรวจสอบจนยืนยันคะแนนครบ
    ซึ่งเป็นเวลาที่บุคลากรต้องใช้จริงเมื่อเก็บข้อมูลด้วยระบบ
\end{enumerate}